\summarychapter{23}{Judgment in the Process of Time}{Wyclif and Hus}

\begin{inbrief}
The Avignon exile and the Great Schism had exposed deep flaws in the papal church, yet after the
councils failed no real reform came from within. Two men, the Englishman John Wyclif and the Czech
John Hus, dared to think of the church as something other than a visible organization headed by the
pope. Both were condemned, and Hus was burned, but they pointed the way to the Reformation.
\question{Who were Wyclif and Hus, and how did they point the way to the future?}
\end{inbrief}

Herbert Butterfield observed in \emph{Christianity and History} that human systems rise, thrive and
fall because time itself carries a built-in judgment: institutions that look worthy crumble as the
centuries bring out their flaws. The idea that the papacy is the channel of God's will died slowly,
for the pope seemed essential both to religion and to the sanctioning of political rule, the center
of Christian society. Wyclif and Hus saw that the hour had come for judgment to fall on the house of
God.

\ssection{The English Zealot, Wyclif}
Wyclif was a zealot who despised neutrality, and he has been praised to heaven or consigned to
darkness ever since. Little is known of the man, who buried his character under pages of scholarly
argument. Raised in northern England, he took his doctorate at Oxford in 1372 and at once became its
leading professor. The burning question was dominion, the right to rule. All agreed that it comes
from God. Some held that it is just only when it comes through the Roman church, to whose pope God
gave universal dominion; others, that it depends on the ruler's being in a state of grace. Wyclif's
teacher Richard FitzRalph asked why grace should be demanded of lay rulers only: do churchmen have
the right to rule while living in mortal sin? Wyclif added that the English government has a
God-given duty to correct abuses in the church within its realm, to remove churchmen who persist in
sin and even to seize the property of corrupt officials. The pope condemned this teaching in 1377,
but powerful friends kept the condemnation to threats.

The lasting importance of the doctrine is its link with the Reformation. ``Dominion founded on
grace'' was Wyclif's way of asserting the spiritual freedom of the righteous: ``God gives no lordship
to his servants without first giving himself to them.'' Priest and layman stand equal before God, the
personal relation with God is everything, and character is the only basis of office. A mediating
priesthood and sacrificial masses are no longer essential. In this Wyclif anticipates Luther's
justification by faith alone.

The decisive year was 1378, when the Great Schism began. Wyclif's view of the papacy had been shaped
by apostolic poverty: the ``Bible papacy'' is a poor and humble life of service, shepherding and
preaching, with no room for temporal power, luxury or politics. At first he welcomed the schism as
proof of the office's spiritual bankruptcy; as it dragged on he came to call the pope Antichrist, a
title two popes merely shared. Christ is truth, lived in poverty and refused temporal dominion; the
pope is the principle of falsehood and seeks magnificence and power. He scorned the argument from
Peter's death in Rome, by which a Muslim might rank a ``prelate of Jerusalem'' above the pope. Christ
alone is head of the church.

\ssection{From Reformer to Protestant}
With this step Wyclif passed from an orthodox preacher of reform to a Protestant, to use a later
word. He kept the old division of the church into the triumphant in heaven, the militant on earth and
those ``asleep in purgatory,'' but he defined the church on earth as the whole number of the elect.
So strict was his predestinarianism that not even a pope can know whether he belongs to it; yet each
should hope that he does, and each who is damned is damned by his own guilt. From this invisible
church of the elect he drew practical conclusions. The church knows no papal primacy or hierarchy
and no ``sects'' of monks, friars and priests, and the salvation of the elect cannot depend on
masses, indulgences, penance or other devices of priestcraft. He went on to challenge pardons,
absolutions, pilgrimages, the worship of images, the adoration of saints, the treasury of merits and
the distinction between venial and mortal sins. He kept purgatory and extreme unction, though he
could not find the latter in the Bible; he tolerated images that aid devotion and prayers to saints;
he valued voluntary confession and called compulsory confession ``the bondage of Antichrist.''
Preaching, he said, ``is of more value than the administration of any sacrament.''

His standard was Scripture. No saint, not Augustine or Jerome, is to be accepted except where he
rests on the Bible, and Christ's law ``is best and enough.'' The New Testament is open to the
understanding of simple people in everything needful for salvation, for Christ wrote his laws in
human hearts. Nothing roused such hostility as his attack on transubstantiation. In the summer of 1380
he published twelve arguments that the consecrated bread and wine are effective symbols: Christ is
present sacramentally, not materially, and the aim of the sacrament is his presence in the soul. His
support at Oxford shrank. The chancellor forbade him to lecture, a council under Archbishop William
Courtenay condemned ten of his doctrines as heretical, and by 1382 he was silenced at the university.

By then he had turned to the people. With a few Oxford scholars he translated the Latin Bible into
English and, copying Francis and the friars, sent out ``poor priests'' in russet robes, barefoot and
with staff in hand, depending on their neighbors for food and carrying pages of his Bible, tracts and
sermons. Their enemies called them Lollards, ``mumblers,'' and they became a power in the land; one
alarmed observer claimed that every second man he met was one. Wyclif's legacy was thus carried mostly
by lay preachers of a simple dependence on the Bible, an example, Shelley notes, of the power of
Scripture read in the people's language. His followers were hunted, expelled from Oxford or forced to
recant, but he himself was left in peace in his parish of Lutterworth, where he died in 1384.

\ssection{The Passage to Bohemia}
The movement survived under restriction in England and found a wider field in Bohemia. The marriage
of Anne of Bohemia to Richard II in 1383 linked the two countries, and students moved between Oxford
and Prague. In Bohemia it joined a strong national party led by John Hus. Born of peasants at
Husinetz in southern Bohemia, he took his bachelor's (1394) and master's (1396) degrees at Prague and
taught in the faculty of arts. As rector and preacher of Bethlehem Chapel he came upon Wyclif's
religious writings and adopted his view of the church as an elect company with Christ, not the pope,
as its head. The chapel, founded in 1391 to foster the national faith, had wall paintings contrasting
pope and Christ: the pope on horseback, Christ barefoot; Jesus washing feet, the pope having his
kissed. Hus's fiery sermons in Czech won wide popular support, and students rioted for and against
Wyclif. On the pope's orders Archbishop Zbynek excommunicated Hus, and tumult followed. When Hus
attacked the pope's sale of indulgences to finance a war against Naples, he lost the support of King
Wenceslas; Prague was placed under interdict, and he withdrew to southern Bohemia, where he wrote
his chief work, \emph{On the Church}, drawing heavily on Wyclif. At the urging of the emperor
Sigismund he agreed to appear at the Council of Constance, hoping to present his views. Instead he
found himself before the Inquisition.

\begin{sidebox}{Inquisition}
Inquisition is a procedure from Roman law: a trial before a panel of judges, all of whom question
the witnesses, which lessens the risk of false testimony. In capital cases Roman law required a
confession or two eyewitnesses; where these were lacking and circumstances suggested guilt, the
accused was sometimes tortured to obtain a confession. The church adopted this arrangement.
\end{sidebox}

Its rule was simple. If enough witnesses testified to guilt, the accused had to confess and renounce
his errors or be burned; the reward of confession was life imprisonment. The judges believed the
witnesses and condemned Hus for heresies he had never taught.

\ssection{Away with Heretics}
Hus was ready to submit to the church if shown from Scripture where he was wrong, but he could not
recant what he had never held. For him truth was supreme: he would not depart from it ``for a chapel
full of gold,'' and his chief fear was that liars would say he had slipped back from what he preached.
For eight months he lay in prison at Constance; the letters and prayers of his last month, asking
for ``a fearless heart, a right faith, a firm hope, a perfect love,'' rank among the greatest in
Christian literature. On July 6, 1415, on the way to the stake he passed a bonfire of his books and
laughed. Asked a last time to recant, he answered that the evidence against him was false, that he
had preached only to win people from their sins, and that in the truth of the gospel he would gladly
die.

The Bohemian revolt did not die with him. Its moderates, the Utraquists, asked for communion in
``both'' bread and cup. Its militants, the Taborites, named after their stronghold, fought the Roman
church and the German empire until wars wore them down. An independent church survived, the
\emph{Unitas Fratrum} or Unity of the Brotherhood, a root in dry ground until Luther came. By the end
of the fifteenth century the dreams of Marsilius had vanished, the conciliar reformers were
repudiated and the revolts of Wyclif and Hus crushed. The period proves, Shelley concludes, that
reform of the papal church from within was impossible. The time of judgment had come.

\begin{keyfacts}
\keypeople{John Wyclif; Richard FitzRalph; Archbishop William Courtenay; John Hus; Anne of Bohemia and
Richard II; Archbishop Zbynek; King Wenceslas; the emperor Sigismund.}
\keydates{%
  \kd{1372}{Wyclif takes his doctorate at Oxford}
  \kd{1377}{The pope condemns Wyclif's teaching on dominion}
  \kd{1380}{Wyclif's twelve arguments against transubstantiation}
  \kd{1382}{Wyclif silenced at Oxford}
  \kd{1384}{Death of Wyclif at Lutterworth}
  \kd{1415}{John Hus burned at Constance (July 6)}}
\keyterms{Dominion founded on grace; invisible church of the elect; Antichrist; authority of
Scripture; Lollards (``poor priests''); Bethlehem Chapel; Inquisition; Utraquists; Taborites;
\emph{Unitas Fratrum}.}
\end{keyfacts}
